% =========================================================
% ppt/src/listing-ppt.tex
% Gaya penampil kode untuk slide.
%
% Diturunkan dari buku_ajar/config/listings.tex, tetapi
% disesuaikan untuk slide: ukuran lebih kecil (karena frame
% 16:9 dan kode harus utuh terlihat) dan jarak lebih rapat.
% =========================================================

% ---------------------------------------------------------
% 1. Definisi bahasa MASM/TASM 8086 (salinan dari buku)
% ---------------------------------------------------------
\lstdefinelanguage{MASM}{
    morekeywords={
        MOV, XCHG, LEA, LDS, LES, PUSH, POP, PUSHA, POPA, PUSHF, POPF,
        LAHF, SAHF, XLAT, IN, OUT, CBW, CWD,
        ADD, ADC, SUB, SBB, INC, DEC, NEG, CMP, MUL, IMUL, DIV, IDIV,
        AAA, AAS, AAM, AAD, DAA, DAS,
        AND, OR, XOR, NOT, TEST, SHL, SHR, SAL, SAR, ROL, ROR, RCL, RCR,
        JMP, JE, JZ, JNE, JNZ, JG, JGE, JL, JLE, JA, JAE, JB, JBE,
        JC, JNC, JS, JNS, JO, JNO, JP, JPE, JNP, JPO, JCXZ,
        LOOP, LOOPE, LOOPZ, LOOPNE, LOOPNZ, CALL, RET, INT, IRET, INTO,
        HLT, NOP, CLI, STI, CLD, STD, CLC, STC, CMC, WAIT, ESC, LOCK,
        DB, DW, DD, DQ, DT, EQU, ORG, SEGMENT, ENDS, ASSUME, END,
        PROC, ENDP, MACRO, ENDM, LOCAL, PUBLIC, EXTRN, EXTERN,
        INCLUDE, INCLUDELIB, MODEL, STACK, CODE, DATA, CONST, ALIGN, EVEN,
        OFFSET, PTR, BYTE, WORD, DWORD, QWORD, TBYTE,
        NEAR, FAR, SHORT, LABEL, TYPE, SIZE, LENGTH, DUP, THIS,
        AX, BX, CX, DX, AL, AH, BL, BH, CL, CH, DL, DH,
        SP, BP, SI, DI, IP, CS, DS, SS, ES, FLAGS,
        CF, PF, AF, ZF, TF, IF, DF, OF
    },
    sensitive=false,
    morecomment=[l]{;},
    morestring=[b]',
    morestring=[b]"
}

% ---------------------------------------------------------
% 2. Gaya kode assembly untuk slide
% ---------------------------------------------------------
\lstdefinestyle{asmslide}{
    language=MASM,
    basicstyle=\ttfamily\scriptsize,
    keywordstyle=\color{obeBlue}\bfseries,
    commentstyle=\color{obeGreen}\itshape,
    stringstyle=\color{obeMaroon},
    numberstyle=\tiny\color{obeGray},
    numbers=left,
    numbersep=5pt,
    stepnumber=1,
    backgroundcolor=\color{codeBack},
    frame=single,
    rulecolor=\color{codeFrame},
    framesep=4pt,
    breaklines=true,
    breakatwhitespace=false,
    % CATATAN PENTING: opsi `postbreak` TIDAK boleh dipakai di sini.
    % Di beamer, `breaklines=true` bersama `postbreak=<kode>` membuat
    % listings gagal dirakit dengan galat
    %   "Missing number, treated as zero <to be read again> \kern".
    % Karena \lstset{style=asmslide} bersifat global, opsi postbreak
    % juga bocor ke gaya lain (mis. outslide/keluaran) dan merusaknya.
    % Di buku ajar (kelas `book`) justru sebaliknya: postbreak aman dan
    % dipakai. Jadi larangan ini khusus untuk slide.
    showstringspaces=false,
    tabsize=4,
    columns=fullflexible,
    keepspaces=true,
    upquote=true,
    xleftmargin=12pt,
    aboveskip=0.4em,
    belowskip=0.4em
}

% ---------------------------------------------------------
% 3. Gaya keluaran program / sesi debugger untuk slide
% ---------------------------------------------------------
\lstdefinestyle{outslide}{
    basicstyle=\ttfamily\scriptsize,
    keywordstyle=\color{black},
    commentstyle=\color{obeGray},
    stringstyle=\color{black},
    numbers=none,
    backgroundcolor=\color{obeGrayLight},
    frame=single,
    rulecolor=\color{codeFrame},
    framesep=4pt,
    breaklines=true,
    showstringspaces=false,
    columns=fullflexible,
    keepspaces=true,
    upquote=true,
    xleftmargin=4pt,
    aboveskip=0.4em,
    belowskip=0.4em
}

\lstset{style=asmslide}

% ---------------------------------------------------------
% 4. Environment ringkas untuk dipakai di dalam frame
%
% Frame yang memuat kedua environment ini WAJIB memakai
% opsi [fragile], karena isinya verbatim.
% ---------------------------------------------------------
\lstnewenvironment{kodeasm}[1][]{%
    \lstset{style=asmslide, #1}%
}{}

\lstnewenvironment{keluaran}[1][]{%
    \lstset{style=outslide, #1}%
}{}
